\section{Introduction}
\label{sec:introduction}

Digital Twins are increasingly presented as foundations for monitoring, simulation, prediction, optimisation, and decision support in manufacturing. Foundational and review literature emphasises virtual--physical linkage, lifecycle information, data integration, simulation, and synchronisation \cite{grieves2017digital,kritzinger2018digital,fuller2020digital,negri2017review,semeraro2021digital}. These perspectives are necessary, but they explain Digital Twin development mainly through architectures, data flows, and model integration. They say less about how operator knowledge is elicited, interpreted, validated, and represented before formal modelling begins.

This omission matters for human-\allowbreak in-\allowbreak the-\allowbreak loop Digital Twins. Operators do not only use dashboards or respond to alarms. They interpret machine behaviour through situated cues such as sound, vibration, rhythm, surface finish, tool behaviour, alarm combinations, and deviations from expected response. Their actions depend on context, uncertainty, production risk, escalation paths, and trust in the available evidence. Human-centred Digital Twin research has therefore begun to ask where humans fit in the Digital Twin loop and how collaboration, responsibility, and trust should be distributed between human actors and digital systems \cite{agrawal2023humans,longo2022ubiquitous,lee2004trust,parasuraman2000levels}. A practical gap remains between recognising that operator knowledge matters and establishing what situated elicitation captures, how researcher interpretation changes through operator review, and whether the validated knowledge can be represented coherently for later engineering use.

We address that gap through the notion of an \emph{operator-strategy fragment}. An operator-strategy fragment is an intermediate knowledge object that relates an operational observation to its context, interpretation, possible responses, rationale, uncertainty, risk, responsibility, and validation state. It is more structured than a raw field note, but less formal than an executable rule or complete Digital Twin model. This intermediate status is important because premature formalisation can erase precisely the contextual judgement that makes operator knowledge useful.

This paper presents \emph{Notebook-to-OSL}. The method is designed for apprenticeship-like situations where a researcher or engineer stands beside an operator and is trained in the work process. During the training session, the researcher captures low-friction digital field notes rather than forcing each observation into a schema. A short post-session interview over the same task scope provides a comparison source. The researcher then annotates the captured material, records an initial interpretation, reviews it with the operator, and transforms validated fragments into candidate formal representations. OSL, short for Operator Strategy Language, is treated as one downstream target. The contribution is not the complete language or a runtime Digital Twin; it is a study design and transformation method that makes the path from situated explanation to formal representation inspectable.

The method is grounded in five bodies of literature. Contextual inquiry supports learning from real work in context and treating interpretation as a collaborative activity \cite{beyer1995apprenticing,beyer1997contextual}. Situated learning and cognitive apprenticeship justify the training situation as an elicitation setting in which expert reasoning becomes visible through demonstration, coaching, articulation, and reflection \cite{lave1991situated,collins1989cognitive}. Cognitive task analysis supports the focus on cues, difficult decisions, hidden judgments, alternatives, and expert--novice differences \cite{crandall2006workingminds,klein1989critical,stanik2021acta}. Requirements engineering and systems engineering provide traceability, rationale, risk, source, status, and validation principles \cite{vanlamsweerde2009requirements,sebokStakeholder,sebokRequirements}. Finally, MBSE and semantic Digital Twin research motivate the assessment of whether validated fragments can support formal representations such as SysML v2, requirements, knowledge graphs, or OSL constructs \cite{omgSysml2024,karabulut2023ontologies,listl2024knowledge,braun2022graph}.

The paper addresses three research questions:

\begin{description}[leftmargin=2.4cm,style=nextline]
    \item[RQ1] How does situated CNC operator training differ from a conventional post-session interview in the operator-strategy knowledge it elicits?
    \item[RQ2] How does operator validation change the researcher's interpretation of captured notes?
    \item[RQ3] To what extent can validated operator-strategy fragments be transformed into traceable, semantically coherent OSL representations, and what gaps remain?
\end{description}

The paper is organised around these questions. For RQ1, it defines a comparative elicitation design that separates material captured during situated training from material elicited in a short post-session interview. For RQ2, it defines a traceable interpretation and validation workflow that records whether fragments are accepted, corrected, narrowed, rejected, or marked sensitive. For RQ3, it defines a model-readiness assessment that tests whether validated fragments can be transformed without losing context, uncertainty, alternatives, provenance, or unresolved knowledge gaps. The annotation schema and illustrative CNC transformation support these three analyses rather than constituting separate contributions detached from the research questions.

The remainder of the paper follows the same structure. Section~\ref{sec:background} establishes the theoretical basis for the three questions. Section~\ref{sec:research-design} maps each question to its data, comparison, and analysis. Section~\ref{sec:method} defines the situated capture, interview, interpretation, validation, and transformation workflow. Section~\ref{sec:schema} specifies the information required to conduct the analyses. Section~\ref{sec:case} illustrates the trace from raw note to formal candidate. Section~\ref{sec:discussion} discusses the contribution and evaluation logic question by question. Section~\ref{sec:conclusion} summarises what the proposed study can establish and what requires empirical application.
